\section{Whole symmetric projective bundles for actual bases}
\label{sec:symmetric-polynomial-basis-comparison}

\begin{theorem}[Actual basis comparison]
\label{thm:symmetric-polynomial-basis-comparison}
\lean{EnvelopeProperty.SymmetricPolynomialBasisComparison.toPolynomial, EnvelopeProperty.SymmetricPolynomialBasisComparison.toSymmetric, EnvelopeProperty.SymmetricPolynomialBasisComparison.iso, EnvelopeProperty.SymmetricPolynomialBasisComparison.iso_structure}
\leanok
\uses{thm:fc-symmetric-projective-embedding}
Let $M$ be a module over any commutative ring $R$, with an actual
basis indexed by a set $I$. Then the entire native projective bundle
$\operatorname{Proj}(\operatorname{Sym}_R M)$ is canonically isomorphic
to the entire projective spectrum of the original polynomial ring
$R[x_i:i\in I]$. The isomorphism preserves the entire original
structural morphism to $\operatorname{Spec}R$.
\end{theorem}

\begin{proof}
The actual algebra isomorphism sends the original basis vectors to
the original polynomial variables and its inverse sends those variables
to the original symmetric generators. Both maps preserve every
homogeneous degree. For the inverse this is the polynomial presentation
lemma. For the forward map, expand each original vector using the actual
basis: its image is homogeneous of degree one. The degree-$d$ symmetric
component is the $d$th power of the degree-one submodule, so its image
lies in the actual degree-$d$ polynomial component. Both original maps
are surjective, and hence satisfy the true irrelevant-ideal condition.
Their two genuine native projective-spectrum maps are inverse by
functoriality and the two original algebra inverse identities. Finally
the original polynomial presentation comparison already retains the
entire coefficient projection. Apply this equality to the inverse
isomorphism and cancel that actual isomorphism. No integrality,
Noetherianity, finite presentation, or supplied projective isomorphism
is required.
\end{proof}
